\section{Phân công công việc}
\label{sec:workload}

Đề bài Project 3 quy định mỗi thành viên tự thực hiện phần đã phụ trách ở Project 2 và mỗi người bổ sung ít nhất hai kiểm thử phi chức năng. Bảng dưới đây tổng hợp vai trò theo phân công Project 2 và cấu trúc thư mục/mã của Project 3.\cite{prj3brief,prj2report,prj3source}

\begin{table}[H]
\centering
\caption{Vai trò thành viên trong Project 3}
\label{tab:team-workload}
\small
\begin{tabular}{|>{\raggedright\arraybackslash}p{3.2cm}|>{\raggedright\arraybackslash}p{4.1cm}|>{\raggedright\arraybackslash}p{5.0cm}|>{\raggedright\arraybackslash}p{2.0cm}|}
\hline
\textbf{Thành viên} & \textbf{Chức năng} & \textbf{Kiểm thử phi chức năng} & \textbf{Hoàn thành} \\
\hline
Nguyễn Quốc Khải & Add a Category & Hiệu năng; tương thích. & 100\% \\
\hline
Đặng Ngọc Khiêu & Create a Group in Course & Hiệu năng; truy cập bằng bàn phím. & 100\% \\
\hline
Đặng Đình Thiên Quân & Add a new user & Hiệu năng; tương thích. & 100\% \\
\hline
Nguyễn Minh Phúc & Create Quiz & Hiệu năng; tương thích. & 100\% \\
\hline
Hồ Tây Nguyên & Add a New Course & Hiệu năng; bảo mật đầu vào. & 100\% \\
\hline
\end{tabular}
\end{table}

Các công việc chung gồm chuẩn bị Moodle và dữ liệu fixture, cấu hình URL/locator, kiểm tra dữ liệu đầu vào, thực thi và lưu kết quả, dọn đối tượng thử nghiệm, tổng hợp kết quả và rà soát sai khác giữa Level 1, Level 2 và môi trường Moodle.
